\section{Polarity, cone reduction, and the scalar input}\label{sec:01-cones}

We first express the volume product as a product of Laplace integrals on
dual cones. This relation is the one used by Klartag
\cite[Lemma~2.1 and Corollary~4.1]{Klartag2018}, with positive duality and
a normalization adapted to the Gaussian argument. We include the
calculation to fix the signs and factorials. The rest of the proof
estimates the integrals after a suitable linear change of coordinates.

\subsection{The center of polarity}
For completeness we establish the minimizing-center convention used in
Theorem~\ref{thm:mahler}. The support function of a convex body $K$ is
$h_K(u)=\sup_{x\in K}\langle u,x\rangle$. Polar coordinates give, for
$z\in\operatorname{int}K$ and $n\ge2$,
\[
 |(K-z)^\circ|=\frac1n\int_{S^{n-1}}
      (h_K(u)-\langle z,u\rangle)^{-n}\,d\sigma(u),
\]
where $\sigma$ is surface measure on the unit sphere.
This is strictly convex as a function of $z$: the function $t\mapsto t^{-n}$
is strictly convex for $t>0$, and for any two distinct centers their
inner products differ on a set of positive surface measure.
It also tends to infinity as $z$ approaches the boundary of $K$.
Indeed, at a limiting boundary point choose a unit supporting normal $u_0$.
Then $\varepsilon=h_K(u_0)-\langle z,u_0\rangle\to0$.
The support functions $h_K(u)-\langle z,u\rangle$ have a common Lipschitz
constant for $z\in K$. On a spherical cap of radius $\varepsilon$ about
$u_0$, their values are at most a constant times $\varepsilon$.
The cap has measure at least a constant times $\varepsilon^{n-1}$,
so its contribution to the integral is at least a constant times
$\varepsilon^{-1}$. Existence and uniqueness of an interior minimizer
follow. In dimension one the same conclusion follows directly from
$(z-a)^{-1}+(b-z)^{-1}$ for $K=[a,b]$.

For an invertible affine map $x\mapsto Bx+a$,
\[
 (BK+a-(Bz+a))^\circ=B^{-\mathsf T}(K-z)^\circ.
\]
The two volume factors change by reciprocal determinants. This proves
affine invariance of the volume product and the equivariance
$s(BK+a)=Bs(K)+a$.

\subsection{The cone formulation}
Put \(m=n+1\). For \(z_0\in\operatorname{int}K\), define
\[
C=\{(y,t):t\ge0,\ y\in t(K-z_0)\},\qquad
D=C^*=\{w:\langle w,x\rangle\ge0\text{ for all }x\in C\}.
\]
For cones, the star denotes this \emph{positive} duality. Both cones are closed, convex, solid, and pointed: solid means having nonempty interior, and pointed means containing no line. For interior vectors \(V\in D\) and \(U\in C\), write
\[
 \chi_C(V)=\int_C e^{-\langle V,x\rangle}\,dx,\qquad
 \chi_D(U)=\int_D e^{-\langle U,x\rangle}\,dx
\]
using \(m\)-dimensional Lebesgue measure.

\begin{lemma}[Cone reduction]\label{lem:cone-reduction}
For \(U=(0,m)\) and \(V=(0,1)\),
\[
\chi_C(V)\chi_D(U)=\frac{(n!)^2}{m^m}|K|\,|(K-z_0)^\circ|.
\]
For any closed convex solid pointed cone and any interior pair \(U\in C\), \(V\in C^*\), the two Laplace integrals are finite. Their product and the inner product \(\langle U,V\rangle\) are preserved by
\[
(C,D,U,V)\longmapsto(BC,B^{-\mathsf T}D,BU,B^{-\mathsf T}V)
\]
for every invertible linear map \(B\).
\end{lemma}
\begin{proof}
A point \((w,t)\) lies in \(D\) exactly when
\(t\ge\sup_{x\in K-z_0}\langle-w,x\rangle\). Since zero is interior to \(K-z_0\), this forces \(t\ge0\), and the slice at height \(t\) is \(-t(K-z_0)^\circ\). Integrating the slices gives
\[
\chi_C((0,1))=n!|K|,\qquad
\chi_D((0,m))=n!m^{-m}|(K-z_0)^\circ|.
\]
For general cone data, interiority of \(V\) gives \(\langle V,x\rangle\ge c|x|\) on \(C\), and interiority of \(U\) gives the analogous estimate on \(D\). These inequalities prove finiteness. Under the displayed transformation, the two change-of-variables factors are \(|\det B|\) and \(|\det B|^{-1}\), and the inner products in the exponential factors are unchanged.
\end{proof}

It therefore suffices to prove
\[
             \chi_C(V)\chi_D(U)\ \ge\ 1,                            \tag{1}\label{eq:1}
\]
whenever \(C\subset\mathbb R^m\) is closed, convex, solid, and pointed, \(D=C^*\), \(U\in\operatorname{int}C\), \(V\in\operatorname{int}D\), and \(\langle U,V\rangle=m\). These will be the standing cone assumptions. We will also show that equality implies that \(C\) is a linear image of an orthant; Section~\ref{sec:06-linear-equality} completes the equality argument for the original body.

The proof constructs two maps from Gaussians using projections on the cones. To control their Jacobians we compare the whole field of projection derivatives to matrix thresholds, with a linear Gaussian matrix as threshold variable. A small feedback in the Gaussian covariance will be used to handle functions of that linear matrix. In the cone argument \(C,D\) refer to cones, while \(\mathsf K,\mathsf C\) below are scalar profiles and \(\mathbf K,\mathbf C\) will be expectations of their symmetric-matrix evaluations. Real layer indices such as \(z\) in the projection field are separate from the Gaussian vector input.

\subsection{Scalar profiles and their identities}
The estimates will use fixed scalar functions, independent of the cone and its dimension. Write \(\phi(x)=(2\pi)^{-1/2}e^{-x^2/2}\), \(p(x)=\int_{-\infty}^x\phi(y)\,dy\), and \(\gamma f=\int f(x)\phi(x)\,dx\). Put \(\mathcal N=x\partial_x-\partial_x^2\). To define three profiles, temporarily write \(X=p/\phi\), \(Y=(1-p)/\phi\), \(f=-(1-p)-\log p\), and \(j=-p-\log(1-p)\). Set
\[
 d=(1+xX)^2 f+(1-xY)^2 j,\qquad
 g=\tfrac12(1-X^2 f-Y^2 j),\qquad v=\log(XY).
\]
\begin{lemma}[Profile identities]\label{lem:profile-identities}
The functions just defined satisfy
\[
 (\mathcal N+2)g=d,\qquad d'=2xd-v'-2g',\qquad
 \mathcal N d-v''=-(\mathcal N+2)(d+g'').                          \tag{2}\label{eq:2}
\]
\end{lemma}
\begin{proof}
We have \(X'=1+xX\), \(Y'=xY-1\), \(f'=-\phi(1-p)/p\), and \(j'=\phi p/(1-p)\). Direct substitution gives
\[
X^2f'+Y^2j'=0,\qquad XX'f'+YY'j'=-1,\qquad
(X')^2f'+(Y')^2j'=-v'.
\]
Also \(X''=xX'+X\) and \(Y''=xY'+Y\). Thus
\(g'=-(XX'f+YY'j)\), and differentiating once more proves \((\mathcal N+2)g=d\). Differentiating \(d=(X')^2f+(Y')^2j\) proves its stated first-order identity. Finally, differentiating the first identity twice gives \(d''=(\mathcal N+4)g''\); combining this with the derivative of the second identity proves the third.
\end{proof}

\subsection{The quantitative scalar inequalities}

The parameters used are (finite decimals throughout denote exact numbers)
\[
\begin{gathered}
 b=.602,\quad c=.365,\quad k=\sqrt{b-c},\quad r=.22,\quad w=4.6,\quad s=3.6,\\
 \eta=.022,\quad \kappa=1.16,\quad \lambda=.65,\quad a_R=7.5,\\
 t_-=-.022,\quad t_+=.064,\quad t_c=.021,\quad t_r=.043,\quad m_*=.352.
\end{gathered}
\]
Define
\[
\begin{aligned}
 \mathsf K&=d-2g,\qquad \mathsf C=-d+(a_R-1)g'',\\
 \pi(x)&=r(\cos(w\operatorname{arsinh}(x/s)),\sin(w\operatorname{arsinh}(x/s))),\\
 q(x)&=k-\sqrt{b-r^2-d(x)},\qquad u=(\pi,q),\qquad S_u=(2+\mathcal N)u=(S_\pi,S_q),\\
 F&=d+|u|^2=c+2kq.
\end{aligned}
\]
The key algebraic choice is that $d+|u|^2$ is affine in $q$.
It will allow a quadratic profile correction to be estimated through one
component. The circular part of $u$ supplies variation on every
nontrivial segment; Appendix~\ref{sec:08-segments} quantifies that fact.
The parameters are fixed independently of the body and dimension.
The first group of inequalities below permits the covariance choice,
the second controls errors from the projection layers, and the third
controls the spectral interval averages of the linear Gaussian field.
Their constants are not asserted to be optimal. The argument uses
precisely the following inequalities, which we collect as one input. Their proofs appear in Appendices~\ref{sec:07-scalar} and~\ref{sec:08-segments}. A bar over a function with specified endpoints denotes its uniform average on that segment, or its value when the endpoints coincide.

\begin{proposition}[Scalar inequalities]\label{prop:scalar-input}
The profiles and exact parameters above have the following properties.

\noindent\textbullet\quad \(b-r^2-d>0,\ |v'-x|\le .319,\ \mathsf C\le-.341,\ S_q>0\), and
\[
 \lambda t_-^2+t_-\mathsf C+\mathsf K>0,\qquad
 \lambda t_+^2+t_+\mathsf C+\mathsf K<0,\qquad
 2(-.37)\mathsf C-(.37)^2-4\lambda\mathsf K>m_*^2 .
\]
\noindent\textbullet\quad Write
\[
 \begin{aligned}
 b_m&=\frac{(1+t_+)(b-3\eta)}{3(3(b+\eta)-4c)},\qquad
 e_0=.256,\qquad \alpha=2.26,\qquad h_s=.5,\\
 H(z,x)&=4c+4kq(x)+2(k-q(x))S_q(z)-2S_\pi(z)\cdot\pi(x),\qquad
 H_0(z,x)=H(z,x)+v''(x)-2\kappa .
 \end{aligned}
\]
Pointwise,
\[
\begin{aligned}
 H+t_c H_0-\tfrac12t_r(h_s+H_0^2/h_s) >
  2(b+\eta+(b+\eta-\kappa)t_-)+2 b_m S_q(z)S_q(x)+2e_0(2t_r)^2+H_0^2/\alpha,\\
 4c+2k(S_q(z)+S_q(x))-S_q(z)S_q(x)-S_\pi(z)\cdot S_\pi(x) >0 .
\end{aligned}                                                       \tag{3}\label{eq:3}
\]
\noindent\textbullet\quad For endpoints \(x_-<x_+\), use subscripts \(\sigma\in\{-,+\}\) for evaluations, with \(-\sigma\) the opposite endpoint. Then
\[
\begin{aligned}
 e_K+\Gamma&+
 \frac{1}{8\lambda}\sum_\sigma
 (\overline{\mathsf C}-\mathsf C_{-\sigma}-\overline{|u-u_\sigma|^2})^2
 +\frac{t_+}{2}\sum_\sigma|\bar u-u_\sigma|^2
 <\overline{|u|^2}-|\bar u|^2,\\
 e_K&=\overline{\mathsf K}-(\mathsf K_++\mathsf K_-)/2,\\
 \Gamma&=\frac{.80}{m_*^2}
 \big[1.25(\mathsf K_+-\mathsf K_-+t_c(\mathsf C_+-\mathsf C_-))^2+
           5 t_r^2(\mathsf C_+-\mathsf C_-)^2\big].
\end{aligned}                                                       \tag{4}\label{eq:4}
\]
For coincident endpoints the two sides in the inequality are equal.
\end{proposition}

The first set of inequalities permits the covariance choice in Section~\ref{sec:02-projections}. The function \(v\) enters the entropy estimate, where Equation~\eqref{eq:2} expresses its contribution through \(d\) and \(g\). The identity \(d+|u|^2=F\) and the two-variable and segment inequalities are used in the later quadratic and spectral estimates. In particular \(d,g,v,q\) are even, as follows from their definitions.

We will also use that these profiles are smooth and that their derivatives
have at most polynomial growth. The exponentially large factor \(X\) does
not make the profiles grow exponentially: its occurrences can be rewritten
in terms of the Gaussian tail ratio \(Y\). On \(x\ge0\),
\(Y=\int_0^\infty e^{-xy-y^2/2}dy\) is comparable to \((1+x)^{-1}\),
and \(1-xY=O((1+x)^{-2})\). Set
\(h_0(y)=(-\log(1-y)-y)/y^2\), with its smooth value at zero.
At \(y=1-p\), the relevant formulas are
\[
 (1+xX)^2f=(1-p(1-xY))^2h_0(y),\qquad
 X^2f=Y^2p^2h_0(y),\qquad j=-p-\log(Y\phi).
\]
Differentiating the Laplace integral for \(Y\) gives bounds for its
derivatives; these formulas then give polynomial derivative bounds for
\(d,g,v,\mathsf K,\mathsf C\). Reflection supplies the negative half-line.
Moreover, \(d\to1/2\) at both ends. The radicand in \(q\) therefore
tends to \(b-r^2-1/2=.0536>0\). Its positivity in
Proposition~\ref{prop:scalar-input}, together with continuity on compact
intervals, gives a uniform positive lower bound. Differentiation of the
square root is thus legitimate globally and gives the same growth property
for \(q\). The bounds for \(\pi\) follow directly from its explicit
formula, and those for \(F\) follow from \(F=d+|u|^2\).
Appendix~\ref{sec:07-scalar}
provides the quantitative bounds used here. No such growth claim is made
for the temporary factors \(X,Y\) separately.

In the cone argument below, a smooth scalar test always means that the
test and all its derivatives have at most polynomial growth. We use
natural logarithms. Matrix inequalities for symmetric matrices are in
quadratic form order.
